\documentclass{article}
\usepackage{neurips_2024}
\usepackage[utf8]{inputenc}
\usepackage[T1]{fontenc}
\usepackage{hyperref}
\usepackage{url}
\usepackage{booktabs}
\usepackage{multirow}
\usepackage{amsfonts}
\usepackage{amsmath}
\usepackage{amssymb}
\usepackage{nicefrac}
\usepackage{microtype}
\usepackage{xcolor}
\usepackage{listings}

\title{Compile-Time Regulatory Compliance:\\
       GDPR Art. 17 as Type Error}
\author{Anonymous Author(s)\\ Paper under double-blind review}
\begin{document}
\maketitle

\begin{abstract}
Data-protection regulations (GDPR Art.~17, EU AI Act Art.~5 and Art.~10,
HIPAA 45\,CFR §164.514) impose operational requirements on the
post-deployment lifecycle of ML models: erasure on request, containment
of prohibited capabilities, composable de-identification.  Current
practice treats these requirements as \emph{external obligations} to be
satisfied by workflow, with compliance argued in human-readable policy
documents.  We demonstrate that --- on Sounio-compiled models with
sedenion-structured activations --- three of these requirements can
be lifted to \emph{compile-time type errors}, each backed by a
\texttt{native\_decide} Lean~4 theorem.  The contribution is not a
new privacy technique; it is a regulatory-to-algebraic reduction
made auditable.
\end{abstract}

\section{Introduction}

Regulatory compliance for ML systems currently relies on workflow.  A
user submits an erasure request (GDPR Art.~17); the engineering team
retrains the model without the user's data; the compliance team
documents the operational event; a certification statement is produced.
Each step involves human judgement, human-written documentation, and
human-audited evidence.  The \emph{algebraic content} --- the actual
statement being certified, ``this model is indistinguishable from one
that never saw the user's data'' --- is not at the compile-time
boundary.

Paper~C of this programme introduced three Sounio compiler-level
wrapper types that make the algebraic content checkable at compile
time:
\texttt{ExactlyPrivate<T>} (erasure),
\texttt{Editable<T>} (edit locality),
\texttt{CapabilityGated<T>} (capability containment).
Paper~G extended these to a closed six-element calculus.
Paper~F --- this paper --- shows that these compiler primitives are
sufficient to express the three-way regulatory surface
\{GDPR Art.~17, EU AI Act Art.~5 and Art.~10, HIPAA Safe Harbor\} as
a library-level desugaring, and that each resulting compliance predicate
is provable by \texttt{native\_decide}.

\paragraph{Methodology disclosure.}
We tag every claim in this paper with one of two markers.
\textbf{[LEAN]} --- mechanically verified by \texttt{lake build} against
the Lean~4 module \texttt{SounioRegulatory.lean}; every compliance theorem
closes by \texttt{native\_decide} over the concrete 168-class ZD structure
of \texttt{SounioZeroDivisorBridge.lean}, with no \texttt{sorry}, no
Mathlib, and no \texttt{omega}-fall-through.
\textbf{[COMPILER]} --- live behaviour of \texttt{bin/souc}: the indicated
\texttt{.sio} program under \texttt{tests/compile-fail/} emits the specified
error code (\texttt{E201}--\texttt{E207}) with the documented message.
This paper makes no operational claim about any deployed system; it is a
correspondence between legal text, type errors, and machine-checked theorems.

\section{Side-by-side: legal text to type error}

\begin{table}[h]
\centering
\caption{[LEAN] + [COMPILER] Regulatory-to-compiler mapping.  Each row pairs a legal text
         with its Sounio type-level encoding and the Lean theorem that
         discharges the algebraic compliance statement.}
\label{tab:reg-map}
\begin{tabular}{p{0.28\linewidth} p{0.28\linewidth} p{0.35\linewidth}}
\toprule
Legal text & Sounio encoding & Lean theorem \\
\midrule
\textbf{GDPR Art.~17(1).} "The data subject shall have the right to
obtain from the controller the erasure of personal data concerning
him or her without undue delay."
&
\texttt{fn unlearn(u: ExactlyPrivate<UserData>) -> ModelWeights}\par
\texttt{with ZD, Witness \{ ... \}}
&
\texttt{gdpr\_art17\_compliance}\newline
$=$ \texttt{unlearning\_kernel\_exact} \\[0.8em]

\textbf{EU AI Act Art.~5(1)(a).} "Placing on the market, putting into
service or use of an AI system that deploys subliminal techniques
beyond a person's consciousness" is prohibited.
&
\texttt{fn deploy(c: CapabilityGated<Model>) -> Output}\par
\texttt{with ZD, Witness \{ ... \}}
&
\texttt{eu\_aiact\_art5\_compliance}\newline
$=$ \texttt{capability\_removal\_preserves\_complement} \\[0.8em]

\textbf{EU AI Act Art.~10(3).} "Training, validation and testing
data sets shall be subject to appropriate data governance and
management practices."
&
\texttt{fn govern(d: ExactlyPrivate<Dataset>) -> Governed}\par
\texttt{with ZD, Witness \{ ... \}}
&
\texttt{gdpr\_art17\_compliance}\newline
(same reduction as Art.~17) \\[0.8em]

\textbf{HIPAA 45\,CFR §164.514(b)(2).} De-identification Safe Harbor
requires removal of the 18 enumerated identifiers.
&
\texttt{fn compose(a: Composable<Cohort>,\par\ \ \ \ \ \ \ b: Composable<Cohort>)}\par
\texttt{with ZD, Witness \{ ... \}}
&
\texttt{hipaa\_safeharbor\_per\_category} $\land$\par
\texttt{hipaa\_composable\_preserves\_deidentification} \\
\bottomrule
\end{tabular}
\end{table}

\section{Compliance-as-type-error semantics}

\paragraph{Operational reduction.}
A Sounio program that type-checks with (e.g.)
\texttt{ExactlyPrivate<T> with ZD, Witness} at every call site that
manipulates personal data is \emph{automatically} in compliance with
GDPR Art.~17(1)'s operational requirement.  The reduction is:

\begin{enumerate}
\item \emph{Compiler gate.} The compiler refuses compilation if any
      call site uses \texttt{ExactlyPrivate<T>} without the
      \texttt{with ZD} effect declared (error 201) or without
      \texttt{with Witness} for the corresponding
      \texttt{Audited<T>} derivative (error 205).
\item \emph{Library gate.} The stdlib module
      \texttt{stdlib/regulatory/gdpr.sio} provides the canonical
      \texttt{gdpr\_erase\_residual} function which returns exactly the
      quantity certified by \texttt{gdpr\_art17\_compliance} to be zero.
\item \emph{Lean gate.} The theorem \texttt{gdpr\_art17\_compliance}
      in \texttt{formal/lean4/SounioRegulatory.lean} discharges via
      \texttt{native\_decide} by reducing to
      \texttt{unlearning\_kernel\_exact}.
\end{enumerate}

\paragraph{What this replaces.}  Current GDPR tooling (e.g.\ 
unlearning libraries, influence-function retraining) provides
\emph{ε-approximate} deletion.  The operational
guarantee after their use is ``the model's behaviour differs from a
never-trained-on-user baseline by at most ε in distribution''.  The
Sounio compile-time gate gives the strictly stronger guarantee that
the model's weights are \emph{bit-equal} to the never-trained
weights on the user's kernel subspace --- a reduction not previously
available to the ML privacy literature~\cite{neel2021descent,
guo2020certified}.

\section{The four compliance theorems}

\texttt{formal/lean4/SounioRegulatory.lean} proves:

\begin{itemize}
\item \texttt{gdpr\_art17\_compliance} ---
      \texttt{alice\_kernel.all (fun v => isZeroPair primA v)}.
      Certified by re-export of \texttt{unlearning\_kernel\_exact}.
\item \texttt{gdpr\_art17\_bounded\_work} ---
      \texttt{alice\_kernel.length = 4}.  Certifies Art.~17's "without
      undue delay" requirement at the primitive level: O(4) kernel
      basis operations per sample.
\item \texttt{eu\_aiact\_art5\_compliance} ---
      \texttt{validPrims.all (fun u => (validPrims.filter (fun v => u ≠ v
      \&\& ¬(isZeroPair u v))).length == 79)}.  Re-export of
      \texttt{capability\_removal\_preserves\_complement}.
\item \texttt{hipaa\_safeharbor\_per\_category} +
      \texttt{hipaa\_composable\_preserves\_deidentification} ---
      the single-category erasure plus the cross-institution merge
      safety.  Each discharged by \texttt{native\_decide} through the
      Paper~G additions (composition preserves orthogonal complement).
\end{itemize}

The top-level \texttt{regulatory\_quadruple} theorem packages all four
into a single conjunction.

\section{Discussion}

\paragraph{What is being claimed.}
Compiling with the specified wrapper types and effect sets produces
source code whose compliance invariants are provable in Lean~4 by
\texttt{native\_decide}.  Compliance with the \emph{textual} regulation
is a human-level judgement that depends on the mapping from legal text
to algebraic invariant --- which we have done explicitly in
Table~\ref{tab:reg-map}.  The formal claim is about the Lean-level
invariant; the regulatory claim is about how well the Lean-level
invariant captures the textual regulation's intent.

\paragraph{What is not being claimed.}
We do not claim that type-checking replaces the human compliance
officer.  The mapping in Table~\ref{tab:reg-map} is a \emph{proposal}
for a legal-to-algebraic translation, and a regulator may disagree with
any row.  What Sounio provides is the compile-time gate: once a regulator
and engineering team agree on the translation, the implementation is
machine-checkable.

\paragraph{Limitations.}
Art.~17(3) of GDPR carries exceptions (freedom of expression, public
interest) that are not algebraic.  The type-level gate only covers the
erasure operation itself; deciding whether the erasure applies in a
given case (data subject eligibility, exception applicability) is out
of scope.  Similarly, EU AI Act Art.~15 ("accuracy, robustness,
cybersecurity") has no obvious ZD-level translation and is not
covered here.

\paragraph{Rebuild requirement.}
The compile-time gates rely on the compiler changes introduced in
Papers~C and~G; see those papers for the bootstrap rebuild prerequisite.

\section{Conclusion}

Three of the four operational requirements of the GDPR-Art.~17 /
EU AI Act Art.~5,10 / HIPAA Safe Harbor regulatory surface reduce to
algebraic invariants on the sedenion ZD graph.  Each invariant is a
compile-time type check in Sounio backed by a \texttt{native\_decide}
Lean theorem.  Regulatory compliance becomes a type-checker's
responsibility, not a policy team's.

\begin{thebibliography}{99}

\bibitem{neel2021descent}
S.~Neel, A.~Roth, and S.~Sharifi-Malvajerdi.
\newblock Descent-to-delete: gradient-based methods for machine unlearning.
\newblock In \emph{ALT}, 2021.

\bibitem{guo2020certified}
C.~Guo, T.~Goldstein, A.~Hannun, and L.~van~der~Maaten.
\newblock Certified data removal from machine learning models.
\newblock In \emph{ICML}, 2020.

\end{thebibliography}
\end{document}
